\begin{table}[htbp!]
  \centering
  \pgfplotstabletypeset[
  col sep=tab,
  text indicator=", string type,
  header=false,
  zerofill,
  columns={0,1,2,3,4},
  columns/0/.style={column name={Linea}, column type=c},
  columns/1/.style={column name={\( n \)}, column type={S[table-format=1.0]}},
  columns/2/.style={column name={
    \text{\( \phi\) \([\si{\degree}]\)}},
    column type={S[table-format=2.1(1)]},
    precision=1,
    sci,
    sci zerofill,
  },
  columns/3/.style={column name={
    \text{\( \lambda \) \([\si{\nano\metre}]\)}},
    column type={S[table-format=3.0(1)]},
    precision=0,
    sci,
    sci zerofill,
  },
  columns/4/.style={column name={
    \text{\( \lambda_e \) \([\si{\nano\metre}]\)}},
    column type={S[table-format=3.1]},
    precision=1,
    sci,
    sci zerofill,
  },
  every even row/.style={
    before row={\rowcolor[gray]{0.9}}
  },
  every head row/.style={
    before row=\toprule, after row=\midrule},
  every last row/.style={after row=\bottomrule}
  ]{./Data/data-hydrogen.tsv}
  \vspace{1.5pt}
  \caption{Longitudes de onda medidas y de referencia para las líneas visibles de la serie de Balmer.}
  \label{tab:lambdas}
\end{table}
